\chapter{Version control, review, and collaborative development}

\section{Version control as a reasoning tool}
Version control records a sequence of states and the intent behind changes. Git
is a distributed revision control system; its commands let a team compare,
select, combine, and recover work. The commands are not the process. A process
is good when it makes change reviewable, integrates often enough to expose
conflicts, and preserves a trustworthy path from request to released state.

A useful commit is coherent: one idea, a message that explains why, and tests
that support the changed behaviour. Small does not mean trivial. A database
migration and the code that requires it may belong in one change if separating
them would create an invalid intermediate state.

\section{Code review}
Review is a risk-reduction activity, not a search for stylistic agreement. Read
the requirement, then the design, then the diff and its tests. Ask:

\begin{itemize}
  \item Is the behaviour correct for valid, invalid, and repeated inputs?
  \item Does the change preserve data, security, compatibility, and performance
        properties that matter?
  \item Are failures observable and recoverable?
  \item Is the new abstraction justified by a real boundary?
  \item Can a future maintainer understand the reason and remove obsolete code?
\end{itemize}

The Git project's own review guidance emphasizes correctness, clarity, and the
importance of reviewing test coverage, while also noting that its advice is not
binding for every project \cite{git2026review}. This is a good example of a
professional recommendation rather than a law.

\section{Review comments}
Prefer comments that identify risk and invite a decision:

\begin{quote}
The retry happens after the database write. If the response is lost, the caller
can submit again and create a second order. Can we store the idempotency key with
the result, or document why duplicate creation is acceptable here?
\end{quote}

This is better than “please refactor.” It connects a concrete path to a property.
Distinguish blocking correctness or security issues from optional suggestions.
Resolve disagreements by testing assumptions or asking the owner of the
requirement to decide, not by escalating the loudest preference.

\section{Documentation that earns its keep}
Documentation serves different readers:

\noindent\begin{tabularx}{\textwidth}{@{}p{0.24\textwidth}X@{}}
\toprule
Document & Primary question \\
\midrule
README & How do I run, test, and orient myself? \\
API reference & What can I call, and what happens? \\
ADR & Why is the architecture this way? \\
Runbook & What should an operator do in a known situation? \\
Comment & Why is this local code non-obvious? \\
\bottomrule
\end{tabularx}

Keep a document near the code it describes, automate examples where possible,
and delete claims that no longer hold. A short, verified runbook is more useful
than a comprehensive page of guesses.

\section{Branching and integration}
Long-lived branches defer integration risk. Short-lived branches and continuous
integration expose conflicts while the author still remembers the change. This
does not require one branching model. A team may use trunk-based development,
release branches, or a hybrid when release governance demands it. What matters is
that the model matches deployment cadence and that it is possible to test the
state being reviewed.

\section{Psychological safety and accountability}
Review quality depends on the social system. People need permission to report a
risk, question a requirement, and admit uncertainty. Psychological safety does
not mean low standards; it makes high standards discussable. Accountability is
stronger when the system records decisions and evidence instead of assigning
blame for every defect.

\begin{exercise}
Write three review comments on a hypothetical change that adds a cache. One
should concern correctness, one privacy or security, and one operability. Mark
which are blocking and why.
\end{exercise}

\section{Chapter summary}
Use version control to preserve intent and reviewable history. Review behaviour,
risk, tests, and operational consequences rather than personal style. Keep
documentation purpose-specific, integrate frequently enough to expose risk, and
build a social environment where uncertainty can be reported early.
